% Sections 15.1 to 15.6, from lectures/L15/appendix/a_theory.md.

\section{Inledning}
\label{c15:sec:inledning}
Komplexa tal behövs när vi vill lösa ekvationer som $x^2 = -1$. Vi inför den \textbf{imaginära
enheten}~$j$:
\[
  j^2 = -1 \quad \Rightarrow \quad j = \sqrt{-1}
\]
Inom elektroteknik används $j$ i stället för $i$, som är reserverat för ström.

\section{Rektangulär form}
\label{c15:sec:rektangular}
Ett komplext tal $z$ skrivs på \textbf{rektangulär form} som
\[
  z = x + jy,
\]
där
\begin{itemize}
  \item $x = \operatorname{Re}(z)$ är den reella delen,
  \item $y = \operatorname{Im}(z)$ är den imaginära delen.
\end{itemize}
Komplexa tal ritas i det \textbf{komplexa talplanet}, med Re-axeln horisontellt och Im-axeln
vertikalt.

\section{Absolutbelopp och fasvinkel}
\label{c15:sec:fasvinkel}
\textbf{Absolutbeloppet} är längden från origo:
\[
  \abs{z} = \sqrt{x^2 + y^2}
\]
\textbf{Fasvinkeln} är vinkeln mot den positiva Re-axeln:
\[
  \delta = \arctan\left(\frac{y}{x}\right) + \text{kvadrantkorrigering}
\]
Korrigeringen beror på i vilken kvadrant talet ligger:

\begin{narrowtable}{@{}l@{\hspace{2.5em}}c@{\hspace{2.5em}}c@{\hspace{2.5em}}l@{}}
  \thead{Kvadrant} & \thead{$x$} & \thead{$y$} & \thead{Korrigering} \\
  \midrule
  I & $+$ & $+$ & ingen \\
  II & $-$ & $+$ & $+180^{\circ}$ \enspace ($+\pi$) \\
  III & $-$ & $-$ & $+180^{\circ}$ \enspace ($+\pi$) \\
  IV & $+$ & $-$ & ingen (negativ $\delta$) \\
\end{narrowtable}

\section{Polär form}
\label{c15:sec:polar}
På \textbf{polär form} skrivs ett komplext tal med sitt absolutbelopp och sin fasvinkel:
\[
  z = \abs{z}\angle\delta
\]
Omvandlingen från \textbf{rektangulär till polär form} ges av
\[
  \abs{z} = \sqrt{x^2 + y^2}, \qquad \delta = \arctan(y/x) + \text{korrigering},
\]
och omvandlingen från \textbf{polär till rektangulär form} av
\[
  x = \abs{z}\cos\delta, \qquad y = \abs{z}\sin\delta.
\]

\section{Räkneoperationer på rektangulär form}
\label{c15:sec:rakneoperationer}
\textbf{Addition och subtraktion:} real- och imaginärdelarna adderas eller subtraheras var
för sig.
\[
  (x_1 + jy_1) \pm (x_2 + jy_2) = (x_1 \pm x_2) + j(y_1 \pm y_2)
\]
\textbf{Multiplikation:} multiplicera ut och använd att $j^2 = -1$.
\[
  (x_1 + jy_1)(x_2 + jy_2) = (x_1x_2 - y_1y_2) + j(x_1y_2 + x_2y_1)
\]
\textbf{Konjugat:} konjugatet till $z = x + jy$ är $\bar{z} = x - jy$. Det är nyttigt vid
division, eftersom det gör nämnaren reell.

\section{Typexempel}
\label{c15:sec:typexempel}

\begin{exempel}[Komplexa tal i talplanet]
Markera $z_1 = 3 + j4$ och $z_2 = -2 - j3$ i det komplexa talplanet och beräkna absolutbelopp och
fasvinkel.

Talet $z_1 = 3 + j4$ ligger i kvadrant I:
\[
  \abs{z_1} = \sqrt{9 + 16} = 5, \qquad
  \delta_1 = \arctan\left(\frac{4}{3}\right) \approx 53{,}1^{\circ}
\]
På polär form är $z_1 = 5\angle 53{,}1^{\circ}$.

Talet $z_2 = -2 - j3$ ligger i kvadrant III, så $180^{\circ}$ läggs till:
\begin{gather*}
  \abs{z_2} = \sqrt{4 + 9} = \sqrt{13} \approx 3{,}61, \\
  \delta_2 = \arctan\left(\frac{-3}{-2}\right) + 180^{\circ} \approx 56{,}3^{\circ} + 180^{\circ}
    = 236{,}3^{\circ}
\end{gather*}
\end{exempel}

\begin{exempel}[Konvertering från rektangulär till polär form]
Spänningen är $U = 3 + j4\,\text{V}$. Bestäm absolutbelopp och fasvinkel.
\[
  \abs{U} = \sqrt{3^2 + 4^2} = 5\,\text{V}, \qquad
  \delta = \arctan\left(\frac{4}{3}\right) \approx 0{,}927\,\text{rad} \approx 53{,}1^{\circ}
\]
På polär form är $U = 5\angle 0{,}927\,\text{rad}$.
\end{exempel}

\needspace{9\baselineskip}
\begin{exempel}[Konvertering från polär till rektangulär form]
Strömmen är $I = 10\angle\frac{\pi}{4}\,\text{mA}$. Bestäm den rektangulära formen.
\begin{gather*}
  x = 10\cos\left(\frac{\pi}{4}\right) = \frac{10}{\sqrt{2}} \approx 7{,}07\,\text{mA} \\
  y = 10\sin\left(\frac{\pi}{4}\right) = \frac{10}{\sqrt{2}} \approx 7{,}07\,\text{mA} \\
  I = (7{,}07 + j7{,}07)\,\text{mA}
\end{gather*}
\end{exempel}
